% Anonymous foundations editions. XeLaTeX, two passes.
\documentclass[10pt,a4paper]{article}
\usepackage[T1]{fontenc}
\usepackage{mathpazo}
\usepackage[scaled=.92]{helvet} % sans-serif vectoriel (évite les polices bitmap EC)
% babel-french complet si disponible ; sinon repli sur les fichiers ini
\IfFileExists{french.ldf}{\usepackage[french]{babel}}{\usepackage[provide=*,french]{babel}}
\usepackage[margin=18mm,top=22mm,bottom=20mm]{geometry}
\usepackage{amsmath,amssymb,stmaryrd}
\usepackage[table]{xcolor}
\usepackage{enumitem}
\usepackage{array,longtable,booktabs}
\usepackage{needspace}
\usepackage[framemethod=default]{mdframed}
\usepackage{fancyhdr}
\usepackage{listings}

%  :  Interrupteur version professeur  :  :  :  :  :  :  :  :  :  :  : --
\newif\ifprof
\ifdefined\eleve\proffalse\else\proftrue\fi

%  :  Palette : marine, acier, ardoise  :  :  :  :  :  :  :  :  :  :  : -
\definecolor{marine}{HTML}{1F3A5F}
\definecolor{acier}{HTML}{4A6FA5}
\definecolor{ardoise}{HTML}{5B6770}
\definecolor{fondpale}{HTML}{F6F8FA}

\setlist[enumerate,1]{label=\arabic*.,leftmargin=*,itemsep=2pt,topsep=3pt}
\setlist[enumerate,2]{label=\alph*.,leftmargin=*,itemsep=1pt,topsep=2pt}
\setlist[itemize]{leftmargin=*,itemsep=2pt,topsep=3pt}
\setlength{\parindent}{0pt}
\setlength{\emergencystretch}{3em}
\setlength{\tabcolsep}{3pt}
\setlength{\parskip}{3pt}

%  :  Symboles mathématiques  :  :  :  :  :  :  :  :  :  :  :  :  :  : --
\newcommand{\N}{\mathbb{N}}
\newcommand{\Z}{\mathbb{Z}}
\newcommand{\Q}{\mathbb{Q}}
\newcommand{\R}{\mathbb{R}}
\newcommand{\C}{\mathbb{C}}
\newcommand{\PP}{\mathcal{P}}
\newcommand{\non}{\operatorname{non}}
\newcommand{\et}{\mathbin{\text{et}}}
\newcommand{\ou}{\mathbin{\text{ou}}}
\newcommand{\Rea}{\operatorname{Re}}   % partie réelle (évite le conflit avec \Re)
\newcommand{\Ima}{\operatorname{Im}}   % partie imaginaire
\newcommand{\card}{\operatorname{card}}
\newcommand{\dd}{\mathrm{d}}
\newcommand{\ii}{\mathrm{i}}           % unité imaginaire (droite)
\newcommand{\vide}{\varnothing}
\newcommand{\bb}[1]{\llbracket #1\rrbracket} % intervalle d'entiers

%  :  Compteur unique des éléments de cours, remis à zéro par chapitre  : 
\newcounter{elt}[section]
\newcounter{seance}[section]
\mdfdefinestyle{filet}{%
  linecolor=acier,linewidth=2.5pt,topline=false,bottomline=false,rightline=false,
  leftmargin=0pt,rightmargin=0pt,innerleftmargin=9pt,innerrightmargin=6pt,
  innertopmargin=4pt,innerbottommargin=4pt,backgroundcolor=fondpale,
  skipabove=5pt,skipbelow=5pt}
\newmdenv[style=filet]{cadre}
\newmdenv[style=filet,linecolor=marine]{objectif}
\newmdenv[style=filet,linecolor=ardoise,backgroundcolor=white]{vigil}

% Étiquette typée : numéro automatique, argument facultatif = nom
\newcommand{\etiq}[2]{\refstepcounter{elt}\textbf{\textcolor{#1}{#2~\thesection.\theelt}}}
\newcommand{\nom}[1]{\if\relax\detokenize{#1}\relax\else\ (#1)\fi\textbf{.}\ }
\newenvironment{definition}[1][]{\begin{cadre}\etiq{acier}{Définition}\nom{#1}}{\end{cadre}}
\newenvironment{propriete}[1][]{\begin{objectif}\etiq{marine}{Propriété}\nom{#1}\itshape}{\end{objectif}}
\newenvironment{theoreme}[1][]{\begin{objectif}\etiq{marine}{Théorème}\nom{#1}\itshape}{\end{objectif}}
\newenvironment{lemme}[1][]{\begin{objectif}\etiq{marine}{Lemme}\nom{#1}\itshape}{\end{objectif}}
\newenvironment{methode}[1][]{\begin{vigil}\etiq{ardoise}{Méthode}\nom{#1}}{\end{vigil}}
\newenvironment{remarque}[1][]{\par\smallskip\noindent\etiq{ardoise}{Remarque}\nom{#1}}{\par\smallskip}
\newenvironment{exemple}[1][]{\par\smallskip\noindent\etiq{ardoise}{Exemple}\nom{#1}}{\par\smallskip}
\newenvironment{demonstration}[1][]{\par\smallskip\noindent\textbf{\textcolor{ardoise}{Démonstration}}%
  \if\relax\detokenize{#1}\relax\else\ (#1)\fi\textbf{.}\ }{\hfill$\square$\par\smallskip}
%  :  En-tête / pied de page  :  :  :  :  :  :  :  :  :  :  :  :  :  :  : 
\newcommand{\titrecourant}{Fondements du supérieur}
\pagestyle{fancy}
\fancyhf{}
\renewcommand{\headrulewidth}{0.4pt}
\fancyhead[L]{\small\textcolor{ardoise}{\ifprof Fondements : enseignant\else Fondements : étudiant\fi}}
\fancyhead[R]{\small\textcolor{ardoise}{\titrecourant}}
\fancyfoot[C]{\small\thepage}
\setlength{\headheight}{14pt}

%  :  Code Python (listings)  :  :  :  :  :  :  :  :  :  :  :  :  :  :  : -
\lstset{basicstyle=\ttfamily\footnotesize,language=Python,columns=fullflexible,
  keepspaces=true,frame=l,framerule=1.5pt,rulecolor=\color{acier},
  xleftmargin=8pt,aboveskip=4pt,belowskip=4pt,showstringspaces=false,
  literate={é}{{\'e}}1 {è}{{\`e}}1 {à}{{\`a}}1 {²}{{\textsuperscript{2}}}1}

\usepackage[hidelinks]{hyperref} % chargé en dernier

\hypersetup{pdftitle={Fondements mathématiques},pdfauthor={},pdfsubject={Cours et préparation de séances},pdfkeywords={}}
